\section{The archimedean oscillator determinant}
\label{sec:oscillator}

\subsection{An equally spaced spectrum}

The Gaussian Mellin integral defines the real archimedean factor.
An additive operator model for its reciprocal must account for an infinite
sequence of equally spaced zeros. The diagonal oscillator is the simplest
such model. Its eigenvalue product diverges, so the normalization must come
from a convergent spectral sum and its analytic continuation.

Let \(H=\ell^2(\mathbb N_0;\mathbb C)\), where
\(\mathbb N_0=\{0,1,2,\ldots\}\). Its inner product is
\(\langle x,y\rangle=\sum_{r\geq0}x_r\overline{y_r}\), linear in
the first variable, and \(e_r\) is the coordinate basis vector.
For positive real \(b\), define
\[
 B_b e_r=br e_r,\qquad
 \mathcal D(B_b)=
 \left\{x\in H:\sum_{r\geq0}b^2r^2|x_r|^2<\infty\right\}.
\]
For positive real \(c\) and \(\operatorname{Re}s>0\), put
\[
 T_{b,c}(s):\mathcal D(B_b)\longrightarrow H,
 \qquad T_{b,c}(s)=\frac{sI+B_b}{c}.
\]
The domain is independent of \(s\). This is a complex operator pencil.
For real \(s>0\), it is positive and self-adjoint.

To check the domain claims, finite sequences give a dense domain.
Coordinate convergence shows that a limit of graph pairs \((x,B_bx)\)
is again a graph pair, so \(B_b\) is closed.
If \(y\) is in the adjoint domain, testing against \(e_r\) forces
\((B_b^*y)_r=br y_r\). This sequence belongs to \(H\) exactly when
\(y\in\mathcal D(B_b)\). The converse follows from the inner product
formula, proving self-adjointness.
For complex \(s\), the same coordinate computation proves normality of
\(T_{b,c}(s)\). Its inverse has diagonal entries \(c/(s+br)\),
which tend to zero. Truncation in the coordinate basis approximates this
inverse in operator norm by finite-rank operators, so it is compact.
Its range lies in \(\mathcal D(B_b)\), since \(br/(s+br)\) is bounded.
The eigenvalues of the pencil are exactly \((s+br)/c\), each of multiplicity
one. Indeed, away from this discrete set the diagonal inverse of the
shifted pencil is bounded and has range in the same domain.

Every eigenvalue lies in the right half-plane. We use its principal
logarithm, with argument in \((-\pi,\pi)\), to define
\begin{equation}
 Z(w,s)=\sum_{r=0}^{\infty}
 \exp\left(-w\operatorname{Log}\frac{s+br}{c}\right),
 \qquad \operatorname{Re}w>1.                     \label{eq:osc-spectral-sum}
\end{equation}
For real \(s>0\), the logarithms of the eigenvalues tend to \(+\infty\),
so their ordinary product diverges. If \(Z\) continues holomorphically
through \(w=0\), its zeta-regularized determinant is, by definition,
\begin{equation}
 D_{b,c}(s)=\det_\infty T_{b,c}(s)
 :=\exp\bigl(-\partial_w Z(0,s)\bigr).             \label{eq:osc-det-definition}
\end{equation}
For the real Hilbert space \(\ell^2(\mathbb N_0;\mathbb R)\), this
definition uses complexification. Each real eigenline becomes one complex
eigenline, so its multiplicity remains one.

\subsection{Continuation and the gamma normalization}

Write \(U=\{a\in\mathbb C:\operatorname{Re}a>0\}\).
On this half-plane, \(\operatorname{Log}\) denotes the principal logarithm,
which is real on the positive axis. All powers of \(a+x\), for \(x\geq0\),
use this logarithm. For \(\operatorname{Re}w>1\), define
\[
 \zeta_H(w,a)=\sum_{r=0}^{\infty}(a+r)^{-w},\qquad a\in U.
\]
On every compact subset of this domain, the summands are bounded by
\(C(1+r)^{-1-\varepsilon}\) for some \(C,\varepsilon>0\).
Thus the sum is jointly holomorphic in \((w,a)\).

\begin{lemma}[A summation formula with a controlled remainder]
\label{lem:osc-summation}
For real \(t\), let \(\{t\}=t-\lfloor t\rfloor\) and set
\(q(t)=\{t\}(1-\{t\})/2\). For every integer \(N\geq0\),
\begin{align}
 \zeta_H(w,a)
 &=\sum_{r=0}^{N-1}(a+r)^{-w}
   +\frac{(N+a)^{1-w}}{w-1}+\frac12(N+a)^{-w}
   +R_N(w,a),                                      \label{eq:osc-cutoff}\\
 R_N(w,a)
 &=w(w+1)\int_N^\infty q(t)(t+a)^{-w-2}\,dt.       \label{eq:osc-remainder}
\end{align}
The integral is jointly holomorphic for \(\operatorname{Re}w>-1\),
\(a\in U\). Formula~\eqref{eq:osc-cutoff} gives a meromorphic continuation
there, with only a simple pole at \(w=1\), of residue \(1\).
In particular,
\begin{equation}
 \zeta_H(0,a)=\frac12-a,\qquad a\in U.              \label{eq:osc-zeta-zero}
\end{equation}
\end{lemma}

\begin{proof}
On \([n,n+1]\), the polynomial \(q(t)=(t-n)(n+1-t)/2\)
vanishes at both endpoints. Its endpoint derivatives are \(1/2\) and
\(-1/2\), and its second derivative is \(-1\).
For any twice continuously differentiable complex function \(f\),
two integrations by parts give
\[
 \int_n^{n+1}q(t)f''(t)\,dt
 =\frac{f(n)+f(n+1)}2-\int_n^{n+1}f(t)\,dt.
\]
Sum this identity from \(n=N\) to \(n=M-1\). The endpoints give
\[
 \sum_{r=N}^{M-1}f(r)
 =\int_N^M f(t)\,dt+\frac{f(N)-f(M)}2
   +\int_N^M q(t)f''(t)\,dt.
\]
Take \(f(t)=(t+a)^{-w}\) and first assume \(\operatorname{Re}w>1\).
Then \(f''(t)=w(w+1)(t+a)^{-w-2}\).
The terms at \(M\) vanish as \(M\to\infty\), while
\(\int_N^\infty f(t)\,dt=(N+a)^{1-w}/(w-1)\).
This proves~\eqref{eq:osc-cutoff} in the initial domain.

For joint continuation, fix compact sets \(K\subset U\) and
\(W\subset\{\operatorname{Re}w>-1\}\).
Put \(\delta=\min_K\operatorname{Re}a>0\) and
\(\sigma=\min_W\operatorname{Re}w>-1\).
The bounds \(0\leq q\leq1/8\) and
\(t+\delta\leq|t+a|\leq t+\max_K|a|\) imply
\[
 |q(t)(t+a)^{-w-2}|
 \leq C_{K,W}(1+t)^{-\sigma-2}.
\]
Indeed, \(|\arg(t+a)|<\pi/2\), and \(\operatorname{Im}w\) is bounded
on \(W\). The same bound for any fixed parameter derivatives acquires
only powers of \(\log(2+t)\) and additional nonpositive powers of \(1+t\).
These bounds remain integrable. Differentiation under the integral is
therefore valid locally in both variables. They also give
\[
 |R_N(w,a)|\leq C_{K,W}|w(w+1)|(1+N)^{-\sigma-1}.
\]
All other terms in~\eqref{eq:osc-cutoff} are holomorphic except the displayed
pole at \(w=1\). Its residue is \((N+a)^0=1\).
The continued formulas for different \(N\) agree, since they agree for
\(\operatorname{Re}w>1\). At zero, the remainder vanishes, and the
remaining terms give \(N-(N+a)+1/2\).
\end{proof}

The derivative at zero involves logarithms of finite products.
The normalization is that of the Gamma integral already used in the
local real factor:
\begin{equation}
 \Gamma(a)=\int_0^\infty e^{-x}x^{a-1}\,dx,
 \qquad a\in U,                                   \label{eq:osc-gamma-integral}
\end{equation}
where \(x^{a-1}=\exp((a-1)\log x)\) uses the real logarithm.

\begin{lemma}[The gamma logarithm and its reciprocal]
\label{lem:osc-gamma}
The integral~\eqref{eq:osc-gamma-integral} is holomorphic and nonzero on
\(U\). There is a unique holomorphic logarithm \(L\) of \(\Gamma\)
which is real on the positive axis. It satisfies, locally uniformly on \(U\),
\begin{equation}
 L(a)=\lim_{N\to\infty}
 \left\{\log((N-1)!)+a\log N
              -\sum_{r=0}^{N-1}\operatorname{Log}(a+r)\right\}.
 \label{eq:osc-gamma-log-limit}
\end{equation}
The reciprocal of \(\Gamma\) extends to the entire function
\begin{equation}
 E(a)=a e^{\gamma a}\prod_{n=1}^{\infty}(1+a/n)e^{-a/n},
 \qquad
 \gamma=\lim_{N\to\infty}\left(\sum_{n=1}^N\frac1n-\log N\right).
 \label{eq:osc-gamma-product}
\end{equation}
Its zeros are simple and occur exactly at \(0,-1,-2,\ldots\).
Moreover, \(\Gamma(1)=1\) and \(\Gamma(a+1)=a\Gamma(a)\) on \(U\).
\end{lemma}

\begin{proof}
On a compact set with \(0<\delta\leq\operatorname{Re}a\leq A\),
the integral and each fixed derivative have an integrable majorant
\[
 e^{-x}(x^{\delta-1}+x^{A-1})(1+|\log x|^k).
\]
This proves holomorphy. Direct integration gives \(\Gamma(1)=1\).
Integration by parts, with vanishing endpoints, gives the recurrence.

For \(N\geq1\), repeated integration by parts on \([0,1]\) gives
\[
 \int_0^1 t^{a-1}(1-t)^N\,dt
 =\frac{N!}{a(a+1)\cdots(a+N)}.
\]
At each step, integrate \(t^{a+j-1}\) and differentiate the remaining
integer power of \(1-t\). The last integral is \(1/(a+N)\).
The substitution \(x=Nt\) yields
\[
 \frac{N!N^a}{a(a+1)\cdots(a+N)}
 =\int_0^N x^{a-1}(1-x/N)^N\,dx
 \longrightarrow\Gamma(a).
\]
For \(0\leq x<N\), one has \((1-x/N)^N\leq e^{-x}\).
Extend the integrand by zero for \(x\geq N\).
The preceding majorants prove convergence locally uniformly in \(a\),
including any fixed number of derivatives.

Let \(H_N=\sum_{n=1}^N1/n\).
The inequality \(\log(1+1/N)>1/(N+1)\) shows that \(H_N-\log N\)
decreases. Integral comparison gives \(H_N\geq\log(N+1)\), so the
sequence has a finite limit \(\gamma\).
The reciprocals of the finite expressions above are
\[
 aN^{-a}\prod_{n=1}^N(1+a/n)
 =a\exp\bigl(a(H_N-\log N)\bigr)
       \prod_{n=1}^N(1+a/n)e^{-a/n}.
\]
For \(|u|\leq1/2\), the logarithm series gives
\( |\operatorname{Log}(1+u)-u|\leq |u|^2\).
Thus the logarithms of the product tails converge uniformly on each
compact subset of \(\mathbb C\). Each tail is the exponential of a
holomorphic sum and is nonzero. The finite initial factors prove that
\(E\) is entire and identify every zero and its order.
Multiplying the two finite expressions and passing to their locally
uniform limits gives \(\Gamma E=1\) on \(U\).

On \(U\), define explicitly
\[
 L(a)=-\operatorname{Log}a-\gamma a
       -\sum_{n=1}^{\infty}
        \bigl(\operatorname{Log}(1+a/n)-a/n\bigr).
\]
The locally uniform convergence proves holomorphy, and
\(\exp L=1/E=\Gamma\). Every term is real for \(a>0\).
Two holomorphic logarithms differ by a constant in \(2\pi i\mathbb Z\),
so the stated real condition proves uniqueness.
Finally, the expression in braces in~\eqref{eq:osc-gamma-log-limit} equals
\[
 -\operatorname{Log}a+a(\log N-H_{N-1})
 -\sum_{n=1}^{N-1}
        \bigl(\operatorname{Log}(1+a/n)-a/n\bigr).
\]
Here \(H_{N-1}-\log N\to\gamma\), which proves the limit with its
specified branch.
\end{proof}

\begin{lemma}[The factorial constant]
\label{lem:osc-stirling}
As positive integers \(N\) tend to infinity,
\begin{equation}
 \log((N-1)!)
 =\left(N-\frac12\right)\log N-N
      +\frac12\log(2\pi)+o(1).                    \label{eq:osc-stirling}
\end{equation}
Also \(\Gamma(1/2)=\sqrt\pi\).
\end{lemma}

\begin{proof}
Let \(J=\int_{\mathbb R}e^{-u^2/2}\,du\), which is finite and positive.
The product integral in polar coordinates gives
\[
 J^2=\int_{\mathbb R^2}e^{-(u^2+v^2)/2}\,du\,dv
 =2\pi\int_0^\infty e^{-r^2/2}r\,dr=2\pi.
\]
Positivity gives \(J=\sqrt{2\pi}\).
Substitution \(x=u^2\) in~\eqref{eq:osc-gamma-integral}, followed by
rescaling this Gaussian integral, gives \(\Gamma(1/2)=\sqrt\pi\).

The recurrence for \(\Gamma\) gives \(n!=\int_0^\infty x^n e^{-x}\,dx\).
Set \(x=n+\sqrt n\,u\). Then
\begin{equation}
 \frac{n!}{n^n e^{-n}\sqrt n}
 =\int_{-\sqrt n}^{\infty}
   \exp\left(n\left[\log(1+u/\sqrt n)-u/\sqrt n\right]\right)du.
 \label{eq:osc-laplace}
\end{equation}
For each fixed \(u\), the exponent tends to \(-u^2/2\).
We give bounds that justify integrating this limit.
For \(-1<x\leq0\), differentiation shows
\(\log(1+x)-x\leq-x^2/2\).
For \(x\geq0\), integration of \(t/(1+t)\) gives
\[
 x-\log(1+x)=\int_0^x\frac{t}{1+t}\,dt
 \geq\frac{x^2}{2(1+x)}.
\]
Consequently the integrand in~\eqref{eq:osc-laplace} is bounded by
\(e^{-u^2/2}\) on \((-\sqrt n,0]\), and by \(e^{-u^2/4}\) on
\([0,\sqrt n]\). Extend these portions by zero to apply dominated
convergence on the two half-lines.

For the remaining tail, \(\log(1+x)/x\) decreases for \(x>0\).
Its derivative has the sign of
\(x/(1+x)-\log(1+x)<0\), as integration of \(1/(1+t)\) shows.
Thus, for \(x\geq1\),
\(\log(1+x)-x\leq-\kappa x\), where \(\kappa=1-\log2>0\).
The tail in~\eqref{eq:osc-laplace} is at most
\[
 \int_{\sqrt n}^{\infty}e^{-\kappa\sqrt n\,u}\,du
 =\frac{e^{-\kappa n}}{\kappa\sqrt n}\longrightarrow0.
\]
It follows that \(n!/(n^n e^{-n}\sqrt n)\to\sqrt{2\pi}\).
Take real logarithms, then use \((N-1)!=N!/N\), to obtain
\eqref{eq:osc-stirling}.
\end{proof}

\begin{proposition}[The two Hurwitz values]
\label{prop:osc-hurwitz-values}
The continuation in Lemma~\ref{lem:osc-summation} is holomorphic near
\(\{0\}\times U\), and
\begin{equation}
 \zeta_H(0,a)=\frac12-a,\qquad
 \partial_w\zeta_H(0,a)=L(a)-\frac12\log(2\pi).
 \label{eq:osc-hurwitz-values}
\end{equation}
Here \(L=\log\Gamma\) is exactly the holomorphic branch constructed in
Lemma~\ref{lem:osc-gamma}.
\end{proposition}

\begin{proof}
The first identity and joint holomorphy were proved in
Lemma~\ref{lem:osc-summation}. Differentiating its exact formula at zero
gives, with \(A=N+a\),
\begin{align*}
 \partial_w\zeta_H(0,a)
 ={}&-\sum_{r=0}^{N-1}\operatorname{Log}(a+r)
       +(A-\tfrac12)\operatorname{Log}A-A\\
    &+\int_N^\infty\frac{q(t)}{(t+a)^2}\,dt.
\end{align*}
The last integral has absolute value at most \(1/(8(N+\delta))\)
when \(\operatorname{Re}a\geq\delta>0\).
It therefore tends to zero locally uniformly on \(U\).
Use~\eqref{eq:osc-gamma-log-limit} to rewrite the remaining expression
as the convergent approximation to \(L(a)\), plus
\[
 -\log((N-1)!)-a\log N
 +(N+a-\tfrac12)\operatorname{Log}(N+a)-(N+a).
\]
By~\eqref{eq:osc-stirling}, this equals
\[
 -\frac12\log(2\pi)
 +(N+a-\tfrac12)\operatorname{Log}(1+a/N)-a+o(1).
\]
For \(a\) in a fixed compact set,
\(\operatorname{Log}(1+a/N)=a/N+O(N^{-2})\) uniformly.
The last two terms therefore tend to zero locally uniformly.
This proves the derivative identity for complex \(a\in U\) directly.
\end{proof}


\subsection{The determinant and the zero mode}

\begin{theorem}[The normalized oscillator determinant]
\label{thm:osc-det}\label{prop:oscillator-det}
For positive real \(b,c\), the spectral sum~\eqref{eq:osc-spectral-sum}
has a jointly meromorphic continuation to
\(\{\operatorname{Re}w>-1\}\times\{\operatorname{Re}s>0\}\).
Its only pole in \(w\) is at \(1\), with residue \(c/b\).
It is jointly holomorphic near \(w=0\), and
\begin{equation}
 D_{b,c}(s)
 =\left(\frac bc\right)^{1/2-s/b}
       \frac{\sqrt{2\pi}}{\Gamma(s/b)}.
 \label{eq:osc-det-formula}
\end{equation}
The power uses the real logarithm of \(b/c\).
This function extends to an entire function of \(s\), with simple zeros
exactly at \(s=-br\), \(r\in\mathbb N_0\).
\end{theorem}

\begin{proof}
Set \(a=s/b\) and \(k=b/c>0\).
Since positive scaling preserves the principal argument,
\(\operatorname{Log}(k(a+r))=\log k+\operatorname{Log}(a+r)\).
Absolute local uniform convergence gives
\[
 Z(w,s)=k^{-w}\zeta_H(w,a).
\]
Lemma~\ref{lem:osc-summation} proves the joint continuation and its
residue. Proposition~\ref{prop:osc-hurwitz-values} gives
\[
 -\partial_w Z(0,s)
 =\left(\frac12-a\right)\log k-L(a)+\frac12\log(2\pi).
\]
Exponentiation proves~\eqref{eq:osc-det-formula} on \(\operatorname{Re}s>0\).
For all complex \(s\), define its continuation by
\[
 \sqrt{2\pi}\exp\bigl((\tfrac12-s/b)\log k\bigr)E(s/b).
\]
The exponential is entire and nonzero. Lemma~\ref{lem:osc-gamma}
identifies the simple zeros of \(E(s/b)\).
\end{proof}

The continuation in this theorem is a scalar entire continuation.
At \(s=-br\), the operator has a zero eigenvalue.
Across a negative real eigenvalue, the principal logarithm has a cut,
so the theorem does not assert a holomorphic principal-cut spectral family
on the whole \(s\)-plane. The specified right half-plane fixes the entire
scalar function uniquely.

\begin{corollary}[Scale and removal of the lowest mode]
For \(d>0\), scaling the operator gives
\begin{equation}
 \det_\infty(dT_{b,c}(s))
 =d^{1/2-s/b}D_{b,c}(s).                           \label{eq:osc-scale}
\end{equation}
The determinant on the closed span of \(e_1,e_2,\ldots\) is
\(D_{b,c}(s+b)\), and
\begin{equation}
 D_{b,c}(s)=\frac{s}{c}D_{b,c}(s+b),\qquad
 D'_{b,c}(0)=\sqrt{\frac{2\pi}{bc}}.               \label{eq:osc-mode}
\end{equation}
The continued determinant with the lowest mode removed has value
\(D_{b,c}(b)=\sqrt{2\pi c/b}\) at \(s=0\).
\end{corollary}

\begin{proof}
Positive scaling replaces \(Z\) by \(d^{-w}Z\).
Differentiating at zero gives~\eqref{eq:osc-scale}, since
\(Z(0,s)=1/2-s/b\).
Reindexing the tail identifies its spectral zeta function with the one
at \(s+b\). The omitted eigenvalue contributes \((s/c)^{-w}\), whose
derivative at zero is \(-\operatorname{Log}(s/c)\).
This proves the first identity in~\eqref{eq:osc-mode} initially in the
right half-plane and then everywhere by entire continuation.
Since \(E(a)/a\to1\) at zero by~\eqref{eq:osc-gamma-product},
\eqref{eq:osc-det-formula} gives the stated slope.
Setting \(s=b\) in that formula gives the final value.
\end{proof}

More generally, recurrence and \(E(a)/a\to1\) give
\[
 D'_{b,c}(-br)=\frac{\sqrt{2\pi}}{b}
       \left(\frac bc\right)^{r+1/2}(-1)^r r!.
\]
Indeed, \(E(a)=a(a+1)\cdots(a+r)E(a+r+1)\), and \(E(1)=1\).
This also verifies the multiplicity at every zero directly.

\subsection{Three scales and the archimedean comparison}

\begin{example}[Unit spacing and unit scale]
For \(b=c=1\), the eigenvalues are \(s,s+1,s+2,\ldots\), and
\[
 D_{1,1}(s)=\frac{\sqrt{2\pi}}{\Gamma(s)},\qquad
 Z(0,s)=\frac12-s.
\]
The integral values \(\Gamma(1)=1\) and \(\Gamma(1/2)=\sqrt\pi\) give
\[
 D_{1,1}(1)=\sqrt{2\pi},\qquad D_{1,1}(1/2)=\sqrt2.
\]
At \(s=1/2\), the exponent in~\eqref{eq:osc-scale} vanishes, so positive
operator rescaling leaves the determinant unchanged. At \(s=1\),
rescaling by \(d\) multiplies it by \(d^{-1/2}\).
The mode identity reads \(D_{1,1}(s)=sD_{1,1}(s+1)\).
It gives a simple zero at zero with slope \(\sqrt{2\pi}\), while the
tail determinant there equals \(\sqrt{2\pi}\).
\end{example}

\begin{example}[The real archimedean scale]
Use Lebesgue measure \(dx\) on \(\mathbb R\), multiplicative measure
\(d^\times x=dx/|x|\) on \(\mathbb R^\times\), and the test function
\(e^{-\pi x^2}\). For \(\operatorname{Re}s>0\), substitution
\(u=\pi x^2\) gives the absolutely convergent integral
\begin{align}
 \Gamma_{\mathbb R}(s)
 &=\int_{\mathbb R^\times}e^{-\pi x^2}|x|^s\,d^\times x
   =2\int_0^\infty e^{-\pi x^2}x^{s-1}\,dx\notag\\
 &=\pi^{-s/2}\int_0^\infty e^{-u}u^{s/2-1}\,du
   =\pi^{-s/2}\Gamma(s/2).                         \label{eq:osc-real-factor}
\end{align}
The factor \(2\) comes from the two real components, and cancels the
Jacobian factor \(1/2\).
For spacing \(b=2\) and operator scale \(c=2\pi\),
\begin{equation}
 \Delta_\infty(s):=D_{2,2\pi}(s)
 =\pi^{s/2-1/2}\frac{\sqrt{2\pi}}{\Gamma(s/2)}
 =\frac{\sqrt2}{\Gamma_{\mathbb R}(s)}.             \label{eq:osc-archimedean}
\end{equation}
Thus \(\Delta_\infty(1)=\sqrt2\) and
\(\Delta_\infty(2)=\sqrt2\,\pi\).
The entire function has zeros at \(0,-2,-4,\ldots\), with
\(\Delta_\infty'(0)=1/\sqrt2\).
Removing the zero mode instead gives value \(\sqrt2\,\pi\) at zero.
These constants use both the measure and the operator scale just stated.
\end{example}

\begin{corollary}[Completed Euler determinant]
\label{cor:completed-euler}
With the measures of Proposition~\ref{prop:local-tate} and scale $2\pi$,
\[
 \frac1{\Delta_\infty(s)\Delta_1(s)}
 =2^{-1/2}\Lambda(s),\qquad \operatorname{Re}s>1.
\]
Both sides continue as single-valued meromorphic scalar functions.
\end{corollary}
\begin{proof}
Theorem~\ref{thm:euler-fredholm} gives $\Delta_1=1/\zeta$ on the
initial half-plane. Equation~\eqref{eq:osc-archimedean} gives
$\Delta_\infty=\sqrt2/\GammaR$. Multiplying proves the identity;
the scalar continuations prove its meromorphic extension.
\end{proof}
The meromorphic identity changes neither the Schatten domain of the
prime diagonal nor the principal spectral cut of the oscillator.
It fixes the normalization a cohomological realization must recover.

\begin{example}[Prime logarithmic spacing]
Let \(p\) be a prime, \(b=\log p\), and \(c=2\pi\).
The heat operator is the diagonal bounded operator with entries
\(e^{-tbr}\). For \(\operatorname{Re}t>0\), these entries have summable
absolute values. Finite truncation therefore gives its trace as
\[
 \operatorname{Tr}(e^{-tB_b})=\sum_{r\geq0}p^{-tr}
 =\frac1{1-p^{-t}}.
\]
The additive determinant has a different spectral parameter:
\[
 D_{\log p,2\pi}(s)
 =\left(\frac{\log p}{2\pi}\right)^{1/2-s/\log p}
       \frac{\sqrt{2\pi}}{\Gamma(s/\log p)}.
\]
At \(t=1\), the trace is \(p/(p-1)\).
At \(s=\log p\), the determinant is \(2\pi/\sqrt{\log p}\).
Its slope at zero is \(1/\sqrt{\log p}\).
After identifying the two parameter planes by \(t=s\), the reciprocal
heat trace has simple zeros at \(2\pi i\mathbb Z/\log p\).
The additive determinant has simple zeros at \(-\mathbb N_0\log p\).
For example, \(-\log p\) is a zero only of the additive determinant.
Hence a nonzero constant cannot identify the two entire functions.
\end{example}


\subsection{Limits of ordinary spectral traces}

\begin{proposition}[Obstructions to ordinary spectral constructions]
\label{prop:spectral-obstructions}
Let \(H\) be the operator of Proposition~\ref{prop:prime-product}.
For every real \(a>0\), its shifted spectral sum
\[
 \sum_{n\geq1}(a+\log n)^{-w}
\]
has no right half-plane of convergence.  In particular, it does not define
a determinant by the spectral-zeta prescription of
\eqref{eq:osc-det-definition}.
There is also no self-adjoint operator \(C\) on a Hilbert space satisfying
\[
 e^{-tC}\in\mathcal S_1,\qquad
 \operatorname{Tr}(e^{-tC})=\GammaR(t)\quad(t>0).
\]
\end{proposition}

\begin{proof}
For real \(w>0\), the terms with \(e^j\leq n<e^{j+1}\) have sum at
least
\[
 \bigl(\lfloor e^{j+1}\rfloor-\lceil e^j\rceil\bigr)
 (a+j+1)^{-w},
\]
which tends to infinity.  For real \(w\leq0\), the summands do not tend
to zero.  Every right half-plane contains such divergent real parameters.
The spectral zeta function has no initial half-plane to continue.

For the second assertion, boundedness of \(e^{-2C}\) forces the
self-adjoint \(C\) to be bounded below, say by \(E_0\).  Functional
calculus and positivity then give
\[
 \operatorname{Tr}(e^{-(2n+2)C})
 \leq e^{-2E_0}\operatorname{Tr}(e^{-2nC}).
\]
But the gamma recurrence gives
\[
 \frac{\GammaR(2n+2)}{\GammaR(2n)}=\frac n\pi,
\]
which is unbounded.  This contradiction proves the claim.
\end{proof}

The first obstruction persists although \(H\) has compact resolvent.
The second concerns ordinary Hilbert traces of a fixed self-adjoint
operator.  The oscillator determinant in
Proposition~\ref{prop:oscillator-det} supplies a regularized construction
with the required archimedean factor.

